\section{Two independent edge periods}\label{two-independent-edge-periods}

Write \(h=2^s\ell_1\), \(k=2^t\ell_2\) and \(d=n-s-t\geq1\). Put

\[R=\{a^2,b^2,c^2,(ab)^h,(bc)^k,(ac)^2,
[(ab)^4,c],[a,(bc)^4]\}.\]

The exact source presentation is

\[G=\begin{cases}
\langle a,b,c\mid R,[(ab)^2,c]^{2^{(d-1)/2}}\rangle,&d\text{ odd},\\
\langle a,b,c\mid R,[(ab)^2,(bc)^2]^{2^{(d-2)/2}}\rangle,&d\text{ even}.
\end{cases}\tag{2}\]

The square-group results UA.Affine and UA.Square identify the quotient
obtained by adding \((ab)^4=(bc)^4=1\). It is the marked affine square
group modulo the lattice

\[\Lambda=\begin{cases}
2^{(d+1)/2}\mathbb Z^2,&d\text{ odd},\\
2^{d/2}\langle(1,1),(1,-1)\rangle,&d\text{ even}.
\end{cases}\tag{3}\]

Indeed, in the affine coordinates of UA.Affine the two commutator words
\[
 [(ab)^2,c],\qquad [(ab)^2,(bc)^2]
\]
are translations by $(0,-2)$ and $(-2,-2)$, respectively. Their powers in (2), and their conjugates under the full signed permutation group $D_8$, generate exactly the two lattices in (3). Normal closure of each period relation is therefore that lattice; this identifies the presented square quotient, not merely its order.
Call it \(B\). It has order \(2^{d+4}\), type \(\{4,4\}\), IP, and every
Boolean deletion. The least scale in (3) is two, including \(d=1,2\);
there is no unverified small seed. These precise presentations and their
normal forms are proved in [U], not inferred from their
orders.

Apply UA.SK to the adjacent pairs, using the dihedral quotients
\(D_{2h}\to D_8\) and \(D_{2k}\to D_8\). The resulting compatible product
\(H\) has type \(\{h,k\}\), IP and order

\[ |H|=2^{d+4}(h/4)(k/4)=2^n\ell_1\ell_2.\tag{4}\]

The relators in (2) hold in it. For the final relator, the base coordinate
is killed by (3), and every attached coordinate is trivial: in either
dihedral factor, one of the commutator entries is trivial or the relevant
square is the identity. The fourth-power commutators also vanish in
every coordinate. Thus the presented group maps onto \(H\).

To recognize the presentation, use UA.Recognition with supports
$\{0,1\}$ and $\{1,2\}$. The local dihedral relations hold by $R$.
The image of the first local kernel is $\langle(ab)^4\rangle$; it is
normal because $a,b$ invert its generator and $c$ centralizes it by
$R$. Likewise $\langle(bc)^4\rangle$ is normal: $b,c$ invert its
generator and $a$ centralizes it. Killing both kernel words gives
exactly the square presentation and the period relation already
identified with $B$. Together with the preceding coordinate check of
all relators, these are the four hypotheses of UA.Recognition.
Consequently $G=H$ as a marked group; in particular (4) is its exact
order. This proves Theorem 2, including its literal presentation.
Solvability follows from the abelian-by-dihedral square base and the
product of the two cyclic projection kernels.
